% Supplement to solver_theory.pdf: the reduced Kirchhoff-metric integrator.
%
% NOTE. The LaTeX source of solver_theory.pdf is not stored in the repository
% (only the built PDF is tracked), so this new material is provided as a
% self-contained, standalone-compilable supplement written in the same notation
% as the parent document. To fold it into the master when that source is
% available, drop the preamble and \input the section bodies after the existing
% "Coupled time integration" section; the cross-references (Lamb impulse L_i,
% Lambda_i, fluid KE T_f, transport term, analytic added mass M_f) are kept
% consistent so it slots in directly.

\documentclass[11pt]{article}
\usepackage[a4paper,margin=1in]{geometry}
\usepackage{amsmath,amssymb,bm}
\usepackage{booktabs}

\newcommand{\R}{\mathbb{R}}
\newcommand{\SEthree}{\mathrm{SE}(3)}
\newcommand{\seiii}{\mathfrak{se}(3)}
\newcommand{\Adstar}{\mathrm{Ad}^{*}}
\newcommand{\adstar}{\mathrm{ad}^{*}}
\newcommand{\Ma}{\mathsf{M}_{a}}
\newcommand{\Ms}{\mathsf{M}_{s}}
\newcommand{\Mmat}{\mathsf{M}}
\newcommand{\Tmap}{\mathsf{T}}
\newcommand{\vv}{\mathbf{v}}
\newcommand{\ww}{\bm{\omega}}
\newcommand{\LL}{\mathbf{L}}
\newcommand{\LamL}{\bm{\Lambda}}
\newcommand{\PP}{\mathbf{P}}
\newcommand{\PiP}{\bm{\Pi}}
\newcommand{\VV}{\mathbf{V}}
\newcommand{\Om}{\bm{\Omega}}
\newcommand{\zz}{\mathbf{z}}
\newcommand{\muv}{\bm{\mu}}
\newcommand{\Fcfg}{\mathbf{F}}
\newcommand{\Tcfg}{\mathbf{T}}
\newcommand{\half}{\tfrac12}
\newcommand{\hatop}[1]{\widehat{#1}}

\title{Supplement: A reduced Kirchhoff-metric variational integrator
for multiple ellipsoids in potential flow}
\date{}

\begin{document}
\maketitle

\begin{abstract}
This supplement develops an alternative to the impulse scheme of the main
document. Because the fluid in potential flow carries no independent state, the
coupled $N$-body system is exactly a finite-dimensional Lagrangian system whose
kinetic energy is a configuration-dependent quadratic form in the body
velocities---the Kirchhoff (added-mass) metric $\Mmat(q)$. We show how to
assemble $\Mmat(q)$ explicitly and cheaply from the same boundary-element (BEM)
solve used elsewhere, cast the dynamics as the body-frame Kirchhoff equations,
and integrate them with a structure-preserving Lie-group scheme: an
implicit-midpoint update whose momentum is transported by the exact $\SEthree$
coadjoint action. The resulting integrator is \emph{fully determined} (there is
no per-body impulse ambiguity), conserves the total linear and angular momentum
maps to machine precision, keeps the energy bounded with no secular drift, and
is second-order accurate. A predictor form reduces the cost to a small multiple
of the impulse scheme's per-step cost.
\end{abstract}

\section{The fluid as a configuration-space metric}
\label{sec:rm-metric}

Collect the rigid-body configuration as $q=\{\mathbf c_i,R_i\}_{i=1}^N$ with
$R_i=R(q_i)\in\mathrm{SO}(3)$ the rotation carried by the unit quaternion $q_i$,
and stack the lab-frame generalized velocity
\begin{equation}
\zz=(\vv_1,\ww_1,\dots,\vv_N,\ww_N)\in\R^{6N}.
\end{equation}
As emphasised in the main document, eqs.~(2)--(3) contain no time derivative of
$\phi$: the potential---and hence every load---is an instantaneous linear
functional of $\zz$ at fixed geometry. In particular the stacked Lamb impulses
of eqs.~(8)--(9),
\begin{equation}
\bm\ell(q,\zz)=(\LL_1,\LamL_1,\dots,\LL_N,\LamL_N),
\end{equation}
are linear in $\zz$ and vanish at $\zz=\mathbf 0$, so they define a symmetric
\emph{added-mass matrix}
\begin{equation}
\boxed{\;\bm\ell(q,\zz)=-\,\Ma(q)\,\zz,\qquad
\Ma(q)=-\,\frac{\partial\bm\ell}{\partial\zz}\;}
\label{eq:rm-Ma}
\end{equation}
The sign is that of the main document: for a single translating body
$\LL_i=-\mathsf M_{f}\vv_i$, so that $\mathbf p_i=m_i\vv_i-\LL_i$ (eq.~(12))
carries the effective inertia $m_i+\mathsf M_f$. The diagonal blocks of $\Ma$
are the (configuration-dependent) self added masses---reducing, for an isolated
body, to the analytic tensors of \S4---and the off-diagonal blocks are the
body--body hydrodynamic coupling. The fluid kinetic energy of eq.~(10) is the
associated quadratic form,
\begin{equation}
T_f(q,\zz)=\half\,\zz^{\!\top}\Ma(q)\,\zz .
\label{eq:rm-Tf}
\end{equation}
Adding the solid kinetic energy $T_s$ gives the reduced Lagrangian
\begin{equation}
\mathcal L(q,\zz)=\half\,\zz^{\!\top}\Mmat(q)\,\zz,\qquad
\Mmat(q)=\Ms(q)+\Ma(q),
\label{eq:rm-lagrangian}
\end{equation}
with the block-diagonal solid metric
$\Ms=\operatorname{blkdiag}\!\big(m_i\mathsf I_3,\;R_iI_i^{\mathrm{body}}R_i^\top\big)$.
Equation~\eqref{eq:rm-lagrangian} is the Kirchhoff/Lamb reduction: the coupled
motion is geodesic motion in the metric $\Mmat(q)$ on the configuration manifold
$Q=\SEthree^{\,N}$.

\paragraph{Why this matters.}
The generalized momentum $\muv=\Mmat(q)\zz$ is a single global covector, and the
equations of motion involve one well-defined configuration force
$\tfrac12\,\partial_q(\zz^\top\Mmat\zz)$. There is no per-body ambiguity: the
under-determination of the raw multibody impulse exchange discussed in the main
development is precisely the absence of the off-diagonal blocks of $\Ma$ and
their configuration gradient. Forming $\Mmat(q)$ makes the exchange explicit.

\section{Explicit assembly of the added-mass metric}
\label{sec:rm-assembly}

The influence matrix $A$ of eq.~(5) depends only on geometry, not on $\zz$, and
its factorisation is already cached (\S7). Column $j$ of $\Ma$ is obtained by
imposing a unit generalized velocity $\zz=\mathbf e_j$ (unit translation or
rotation of one body about one axis), solving the \emph{same} BEM system for the
resulting $\phi$, and reading the stacked impulse~\eqref{eq:rm-Ma}. All $6N$
right-hand sides share the cached factorisation, so the entire matrix costs one
factorisation and $6N$ back-substitutions/warm-started GMRES solves rather than
$6N$ independent solves. Symmetrising $\tfrac12(\Ma+\Ma^\top)$ removes the small
discretisation asymmetry.

This assembly was verified directly. At a close two-body configuration
($\mathrm{ndiv}=2$) the assembled $\Ma$ is symmetric to $3\times10^{-5}$
(relative), reproduces the solver's fluid kinetic energy through
\eqref{eq:rm-Tf} to $6\times10^{-4}$, and its configuration gradient contracted
with the velocity, $\half\,\zz^\top(\partial_q\Ma)\zz$, matches the
finite-differenced $\partial_q T_f$ to $2\times10^{-5}$. This last identity is
the body--body exchange force that the impulse scheme's pair-metric correction
approximates, here obtained exactly. The full matrix costs about $1.9\times$ a
single impulse solve, confirming the shared-factorisation economics.

\section{Body-frame Kirchhoff equations}
\label{sec:rm-kirchhoff}

Working in the lab frame at fixed $\zz$ and differentiating $T$ with respect to
orientation double-counts the kinematic link between $R_i$ and $\ww_i$, injecting
a spurious $-\ww_i\times\mathbf H_i$ torque for anisotropic bodies. We therefore
pass to \emph{body-frame} velocities and momenta, in which each body's self block
is constant. For body $i$ let
\begin{equation}
\VV_i=R_i^\top\vv_i,\quad \Om_i=R_i^\top\ww_i,\qquad
\zeta=(\VV_1,\Om_1,\dots),\qquad
\Mmat_{\mathrm{bf}}(q)=\Tmap(R)^\top\,\Mmat(q)\,\Tmap(R),
\end{equation}
with $\Tmap(R)=\operatorname{blkdiag}(R_i,R_i)$, and body-frame momenta
$\muv=\Mmat_{\mathrm{bf}}\zeta$, written per body as $(\PP_i,\PiP_i)$. The reduced
Euler--Poincar\'e (Kirchhoff) equations on $\SEthree^{\,N}$ are
\begin{align}
\dot\PP_i &= \PP_i\times\Om_i + \Fcfg_i,\label{eq:rm-Pdot}\\
\dot\PiP_i &= \PiP_i\times\Om_i + \PP_i\times\VV_i + \Tcfg_i,\label{eq:rm-Pidot}
\end{align}
where $(\Fcfg_i,\Tcfg_i)$ is the body-frame configuration force, i.e. the
body-referred gradient of the interaction energy at fixed $\zeta$. For a single
body $\Mmat_{\mathrm{bf}}$ is constant, $(\Fcfg,\Tcfg)=\mathbf 0$, and
\eqref{eq:rm-Pdot}--\eqref{eq:rm-Pidot} reduce to the classical Kirchhoff top;
the gyroscopic terms $\PP\times\Om$ and $\PiP\times\Om+\PP\times\VV$ are exactly
what the naive lab-frame gradient omitted.

\subsection*{Configuration force and the momentum identities}
The interaction energy is invariant under a rigid motion of the whole
configuration at fixed body-frame $\zeta$; hence the true configuration force is
$L^2$-orthogonal to the six global rigid-motion directions of $Q$, which
expresses the discrete Noether identities
\begin{equation}
\sum_i R_i\Fcfg_i=\mathbf 0,\qquad
\sum_i\big(\mathbf c_i\times R_i\Fcfg_i+R_i\Tcfg_i\big)=\mathbf 0 .
\label{eq:rm-noether}
\end{equation}
A finite-difference estimate of $(\Fcfg,\Tcfg)$ violates \eqref{eq:rm-noether}
only at truncation order; projecting it onto the complement of the six
rigid-motion directions restores both identities exactly and is essential for
machine-precision momentum conservation.

\section{Structure-preserving time integration}
\label{sec:rm-integrator}

Write $g=(R,\mathbf c)\in\SEthree$ per body and let $\muv=(\PP,\PiP)$ transform
under the coadjoint action of $f=(R_f,\mathbf p_f)\in\SEthree$ as
\begin{equation}
\Adstar_f\,\muv=\big(R_f^\top\PP,\;R_f^\top(\PiP+\PP\times\mathbf p_f)\big),
\label{eq:rm-adstar}
\end{equation}
whose infinitesimal form is the right-hand side of
\eqref{eq:rm-Pdot}--\eqref{eq:rm-Pidot}. The spatial momentum
$\Adstar_{g^{-1}}\muv$ is the conserved Noether quantity. One implicit-midpoint
step, with unknown midpoint velocity $\zeta_{\mathrm{mid}}$, is
\begin{align}
\zeta_{\mathrm{mid}}
&=\Mmat_{\mathrm{bf}}(q_{\mathrm{mid}})^{-1}
   \Big(\Adstar_{f_{n\to\mathrm{mid}}}\muv_n+\half\,\Delta t\,\mathbf f_{\mathrm{cfg}}\Big),
   \label{eq:rm-fixedpoint}\\
q_{n+1}&=q_n\oplus\Delta t\,\zeta_{\mathrm{mid}}
  \quad(\mathbf c_i{+}{=}\Delta t\,R_{\mathrm{mid},i}\VV_i,\;
        R_i\!\leftarrow\!R_i\exp(\Delta t\,\hatop{\Om}_i)),\label{eq:rm-config}\\
\muv_{n+1}&=\Adstar_{f_{\mathrm{mid}\to n+1}}
   \Big(\Adstar_{f_{n\to\mathrm{mid}}}\muv_n+\Delta t\,\mathbf f_{\mathrm{cfg}}\Big).
   \label{eq:rm-commit}
\end{align}
Here $q_{\mathrm{mid}}$ is the half-configuration, $\mathbf f_{\mathrm{cfg}}$ the
projected configuration force there, and each leg increment
$f=(R_a^\top R_b,\,R_a^\top(\mathbf c_b-\mathbf c_a))$ is the \emph{exact} relative
motion between the two configurations it connects, so the two legs
$n\!\to\!\mathrm{mid}\!\to\!n{+}1$ in \eqref{eq:rm-commit} compose to the exact
full increment $n\!\to\!n{+}1$. Consequently
\begin{equation}
\Adstar_{g_{n+1}^{-1}}\muv_{n+1}
 =\Adstar_{g_n^{-1}}\muv_n+\Delta t\sum_i(\text{spatial }\Fcfg_i),
\end{equation}
and the sum vanishes by \eqref{eq:rm-noether}: \textbf{the total linear and
angular spatial momentum are conserved to machine precision}. Two points make
this exact rather than $O(\Delta t^2)$: (i) the impulse is applied \emph{once},
at the half configuration where the projection is performed, so its
transported net force cancels; splitting it across the two legs would leave a
term in the $n{+}1$ frame and leak momentum at $O(\Delta t^2)$ per step; and
(ii) the leg increments are the true relative motions, not the exponential of an
averaged twist.

Energy is not conserved exactly---as for any momentum-conserving variational
integrator---but is bounded with no secular drift. For the free Kirchhoff top,
\eqref{eq:rm-fixedpoint} with $\mathbf f_{\mathrm{cfg}}=\mathbf 0$ conserves both
the energy $\tfrac12\zeta^\top\Mmat_{\mathrm{bf}}\zeta$ and the Casimir
$|\PiP|$ to machine precision.

\subsection*{Predictor form for efficiency}
Evaluating $\Mmat_{\mathrm{bf}}$ and $\mathbf f_{\mathrm{cfg}}$ inside the
fixed point makes each iterate a full metric assembly. Because the metric at the
true versus predicted midpoint differ by $O(\Delta t^2)$, we instead assemble
both \emph{once} per step at a midpoint predicted from the previous step's
$\zeta_{\mathrm{mid}}$, and take a single direct (BEM-free)
update~\eqref{eq:rm-fixedpoint}--\eqref{eq:rm-commit}. Carrying $\muv$ as the
state removes any endpoint metric solve, and the configuration force uses
forward differences. Crucially the \emph{same} predicted midpoint is used for
the projection and for the transport intermediate in \eqref{eq:rm-commit}, so
momentum remains machine-exact while second order is preserved.

\section{Properties and verification}
\label{sec:rm-results}

Table~\ref{tab:rm-results} summarises the behaviour on a spinning single
ellipsoid and a close two-body pair ($1\!:\!0.8\!:\!0.6$ ellipsoids, sep~$3.5$,
$\mathrm{ndiv}=2$, $\Delta t=0.05$). Momentum is conserved to $\sim10^{-13}$;
energy is bounded; the trajectory is second order (error $\propto\Delta t^2$
against a fine reference). Compared against the production variational scheme on
the same close pair, both integrators converge to the same continuum trajectory,
and the reduced-metric scheme is about an order of magnitude more accurate at a
given step, while its predictor form runs at a small multiple of the impulse
scheme's per-step cost.

\begin{table}[h]
\centering
\begin{tabular}{lccc}
\toprule
case / scheme & energy drift & $|\Delta\mathbf P|$ & $|\Delta\mathbf H|$ \\
\midrule
single ellipsoid, free top & $9\times10^{-14}$ & $2\times10^{-14}$ & $3\times10^{-14}$\\
two ellipsoids, close pair & $8\times10^{-5}$ (bounded) & $9\times10^{-14}$ & $1\times10^{-13}$\\
\midrule
trajectory error vs.\ fine ref. ($\Delta t{=}0.05$) & \multicolumn{3}{c}{$1.1\times10^{-4}$ (2nd order)}\\
production variational, same case & \multicolumn{3}{c}{$5.5\times10^{-3}$ (1st order)}\\
\midrule
cost: impulse / reduced-metric (once-per-step) & \multicolumn{3}{c}{$0.12$ / $0.29$ s per step}\\
\bottomrule
\end{tabular}
\caption{Reduced-metric integrator: conservation, accuracy and cost.
Momentum drifts are the maxima over the run; energy drift is relative.}
\label{tab:rm-results}
\end{table}

\paragraph{Relation to the impulse scheme.}
The impulse scheme evaluates only $\bm\ell(q,\zz)$ (one BEM solve) and differences
it; it never forms $\Ma$ or $\partial_q\Ma$, which is both why it is fast and why
its multibody exchange is under-determined. The reduced-metric scheme pays for
the metric and its gradient and, in return, is fully determined and
momentum-exact by construction. The remaining cost gap is entirely the
finite-difference configuration force (a geometry re-mesh per perturbation);
replacing it with an analytic added-mass shape derivative (Hadamard/adjoint) is
the route to full parity and is the natural next development.

\end{document}
